Battery--supercapacitor hybrid energy storage in electric vehicles needs a fast DC-bus voltage loop and a slower power-sharing policy. A reproduction of a recent disturbance-observer sliding-mode controller (DO-SMC) showed that its fixed switching gain violates the controller's own reaching condition by a factor of about forty during acceleration. The controller then regulates the bus no better than a PI loop with the same small-signal poles.

We propose a two-tier architecture. In tier 1, an adaptive DO-SMC estimates the observer error from the identity $\dot{\hat D}=\lambda\tilde D$ and schedules the switching gain on it, so the reaching condition is enforced online for any observer gain. A boundary layer that scales with the gain keeps the small-signal poles fixed. In tier 2, a TD3 agent adapts the filter time constant, the supercapacitor energy target and the observer gain through a safety filter.

The method was evaluated in simulation on standard drive cycles that were never used for training or checkpoint selection, with three training seeds, paired noise seeds, a non-learning SOC-adaptive baseline and a filter ablation.
\begin{itemize}
\item Tier 1 reduced the maximum bus deviation from \R{pt:dosmc:maxdev}\,V to \R{pt:afix:maxdev}\,V on the reproduced transient. This matches an observer-compensated robust MPC at the update cost of the original controller.
\item The learned supervisor lowered the RMS battery current by up to \R{wh:v2vsafix:dibrmsabs}\,\% on cycles resembling the training data, and the safety filter prevented the bus collapses seen without it.
\item On US06 the agent caused power deficits (\R{us06a:v2:sat:mean}\,\% saturated time, none with fixed parameters), and on steady cruising it tripled the battery current. Diagnosis shows the agent used the supercapacitor reference as a fast power actuator, which a filter on the reference value cannot prevent.
\end{itemize}
The results support the certificate-enforcing inner loop and identify the requirements for a deployable learned supervisor.
